\documentclass[11pt]{article}
\usepackage[utf8]{inputenc}
\usepackage{amsmath,amssymb}
\usepackage[margin=1in]{geometry}
\usepackage{hyperref}
\emergencystretch=3em

\title{The complete expansion for every dimension:\\
$H_m(L)=\sum_{i\le m}c_{m,i}\log^iL+O(1/L)$ and $Z_d(n)=n^{-1/2}P_d(\log n)+O(n^{-1})$}
\author{Claude, with Daniel Murfet}
\date{2026-09-06}

\begin{document}
\maketitle
\sloppy

\begin{abstract}
Units~41 and~42 identified the constant terms below $\log n$ and $(\log n)^2$ by hand. This unit
proves the underlying transfer step once, for the iterated logarithmic primitives
$H_m=\texttt{iterLogPrim}\,m$ of unit~39: with explicit recursively defined coefficients
$c_{0,0}=A$, $c_{m+1,i+1}=c_{m,i}/(i+1)$, $c_{m+1,0}=\int_0^1H_m/x+\int_1^\infty\theta_m/x$,
\[
H_m(L)=\sum_{i=0}^mc_{m,i}\log^iL+\theta_m(L),\qquad|\theta_m(L)|\le\frac ML\quad(L\ge1),
\qquad\theta_{m+1}(L)=-\int_L^\infty\frac{\theta_m(x)}{x}\,dx,
\]
and $c_{m,m}=A/m!$. Hence every all-equal chart integral has a complete expansion
$Z_{k+2}(n)=n^{-1/2}\sum_{i\le k+1}c_{k+1,i}(\tfrac12\log n+(k+2)\log b)^i+O(n^{-1})$.
Zero \texttt{sorry}/\texttt{axiom}.
\end{abstract}

\section{Setting}
The recursion $H_{m+1}(L)=\int_0^LH_m(x)/x\,dx$ preserves the shape ``polynomial in $\log x$ plus a
remainder $O(1/x)$'': over $[1,L]$ the polynomial integrates termwise via
$\int_1^L\log^ix/x=\log^{i+1}L/(i+1)$, the remainder over $[1,\infty)$ converges to a constant that
becomes the new $c_{m+1,0}$, and the tail over $[L,\infty)$ is the new remainder. The remainder bound
$M/L$ is preserved exactly because $\int_L^\infty Mx^{-2}=M/L$.

\section{The things formalised}
{\small
\begin{verbatim}
noncomputable def iterLogCoeff (β a : ℝ) : ℕ → ℕ → ℝ
  | 0, 0 => quadMass β a
  | 0, _ + 1 => 0
  | m + 1, 0 => (∫ x in Ioc 0 1, iterLogPrim β a m x / x)
      + ∫ x in Ioi 1, (iterLogPrim β a m x
          - ∑ i ∈ Finset.range (m + 1), iterLogCoeff β a m i * Real.log x ^ i) / x
  | m + 1, i + 1 => iterLogCoeff β a m i / ((i : ℝ) + 1)
theorem iterLogCoeff_top ... : iterLogCoeff β a m m = quadMass β a / (m.factorial : ℝ)
theorem iterLogRem_succ_eq ... (hL : 1 ≤ L)
    (hb : ∀ L, 1 ≤ L → |iterLogRem β a m L| ≤ M / L) :
    iterLogRem β a (m + 1) L = -∫ x in Ioi L, iterLogRem β a m x / x
theorem iterLogPrim_expansion (β a : ℝ) (hβ : 0 < β) (m : ℕ) (L : ℝ) (hL : 1 ≤ L) :
    |iterLogPrim β a m L
        - ∑ i ∈ Finset.range (m + 1), iterLogCoeff β a m i * Real.log L ^ i|
      ≤ quadMoment β a / L
theorem iterChart_expansion_isBigO (β a b : ℝ) (hβ : 0 < β) (hb : 0 < b) (k : ℕ) :
    (fun n => iterChart β a b (k + 2) (Real.sqrt n)
        - n ^ (-(1/2 : ℝ)) * iterChartCoeff β a b k n) =O[atTop] fun n => n⁻¹
\end{verbatim}
}

\section{Proof strategy}
The coefficients are defined by structural recursion on $m$ with the index $i$ as a varying second
argument, so all three equation lemmas are \texttt{rfl}. The remainder bound is proved by induction
on $m$; the inductive step is the exact identity $\theta_{m+1}(L)=-\int_L^\infty\theta_m/x$, which
takes the current bound as a hypothesis (it is needed for integrability of $\theta_m/x$ on
$[1,\infty)$) and is proved by the same three-way split as at $d=3$: $(0,1]$, $(1,L]$ with the
pointwise decomposition $H_m/x=\sum c_{m,i}\log^ix/x+\theta_m/x$, and $\int_1^\infty=\int_1^L+
\int_L^\infty$. The re-indexing $\sum_{j\le m+1}c_{m+1,j}\log^jL=c_{m+1,0}+\sum_{i\le m}
c_{m,i}\log^{i+1}L/(i+1)$ is \texttt{Finset.sum\_range\_succ'} plus the coefficient equation.
The base case is $|F(L)-A|=A-F(L)\le M/L$ from unit~35. The chart statement is unit~39's reduction
$Z_{k+2}=n^{-1/2}H_{k+1}(\sqrt nb^{k+2})$ with $\log(\sqrt nb^{k+2})=\tfrac12\log n+(k+2)\log b$.

\section{Roadblocks and resolutions}
\begin{itemize}
\item The file compiled without errors on the first pass; the only fixes were deprecation renames
(\texttt{continuousOn\_finsetSum}, \texttt{integrable\_finsetSum}, \texttt{integral\_finsetSum}), an
unused hypothesis, and three over-long lines.
\end{itemize}

\section{What was Mathlib, what was new}
Mathlib: \texttt{integral\_finsetSum}, \texttt{integrable\_finsetSum},
\texttt{continuousOn\_finsetSum}, \texttt{Finset.sum\_range\_succ'},
\texttt{norm\_integral\_le\_of\_norm\_le}, and the interval-to-set integrability bridge
\texttt{intervalIntegrable\_iff\_integrableOn\_Ioc\_of\_le}. New: the recursive coefficient scheme,
the exact remainder recursion, and the uniform $M/L$ bound for every $m$.

\section{Lessons}
Two hand-worked instances (units~41, 42) were enough to see the induction; formalising the third
level as a general step cost less than either instance because the statement had already stabilised.
The remainder bound being \emph{exactly} preserved ($M/L\mapsto M/L$) is what makes the induction
hypothesis clean; a bound that degraded with $m$ would have needed the \texttt{iterConst} bookkeeping
of unit~39.

\section{Outcome}
For every dimension $d\ge2$ the all-equal chart integral has a complete polynomial-in-$\log n$
expansion with explicit recursive constants and $O(n^{-1})$ remainder, subsuming units~41--42. Next
candidates: the weighted primitives $F_\gamma$ for general $(k,h)$, or the mixed-multiplicity case in
general dimension (several minimal coordinates plus several non-minimal ones).
\end{document}
